\documentclass[11pt]{article}
\usepackage{amsmath,amssymb,amsthm}
\newtheorem{theorem}{Theorem}
\newtheorem{lemma}[theorem]{Lemma}
\title{A Reflection-Symmetric Cauchy Sign Lift}
\author{}
\date{}
\begin{document}
\maketitle

\paragraph{Status.}
This is a complete written finite sign-interpolation theorem, with an
exact rational audit of all $2^{12}$ limiting sign tables. It is not a
Lean theorem. The conversion below uses twenty original admissible
heads on a twelve-row grid, so this result does not improve the
worst-case asymptotic upper bound. Its useful feature is a proved
target-dependent cancellation mechanism, not its present head count.

Put
\[
 S=\begin{pmatrix}-1&1&-1&1\\-1&-1&1&1\\-1&-1&-1&-1\end{pmatrix},
 \qquad c=(1,-1,-1,1)^\top.
 \tag{1}
\]
Write $v_h$ for its four columns. The matrix $S$ has rank three and
kernel $\mathbb Rc$. Thus $T=S\operatorname{diag}(c)$ has rank three
and kernel $\mathbb R(1,1,1,1)^\top$.

\begin{theorem}[Eight masked Cauchy columns]
Fix any four distinct positive numbers $d_1,\ldots,d_4$. For every
sign table $f_{a\ell}\in\{-1,1\}$, $1\le a\le4$, $1\le\ell\le3$,
there exist positive distinct numbers $t_1,\ldots,t_4$ and real
coefficient functions $q_h^+(\epsilon),q_h^-(\epsilon)$ such that
\[
 F_{a\ell}(\epsilon)=\sum_{h=1}^4
 \left(\frac{v_{\ell h}q_h^+(\epsilon)}{d_a-\epsilon t_h}
       -\frac{v_{\ell h}q_h^-(\epsilon)}{d_a+\epsilon t_h}\right)
 \tag{2}
\]
has $f_{a\ell}F_{a\ell}(\epsilon)>0$ for all sufficiently small
positive $\epsilon$. All denominators in (2) are then positive.
The poles and coefficients can all be chosen rational if the nodes
$d_a$ are rational.
\end{theorem}

\begin{proof}
We first determine an explicit flat limit. Pairing the two columns
associated with $v_h$ and $-v_h$ gives, by invertible column operations,
\[
 E_h(d,\epsilon)=\frac{dv_h}{d^2-\epsilon^2t_h^2},\qquad
 O_h(d,\epsilon)=\frac{v_h}{d^2-\epsilon^2t_h^2}.
 \tag{3}
\]
Indeed, these are respectively half the difference and
$(2\epsilon t_h)^{-1}$ times the sum of the two columns in (2).
For $u_h=t_h^2$ and $w=Tu$,
\[
 \epsilon^{-2}\sum_h c_hE_h\longrightarrow d^{-3}w,
 \qquad
 \epsilon^{-2}\sum_h c_hO_h\longrightarrow d^{-4}w.
 \tag{4}
\]
The constant terms cancel because $Sc=0$. Since $S$ has rank three,
ordinary linear combinations of (3) independently approach $d^{-1}A$
and $d^{-2}B$ for any $A,B\in\mathbb R^3$. Consequently (2) can
approach every vector-valued function
\[
 d^{-1}A+d^{-2}B+\gamma d^{-3}w+\delta d^{-4}w.
 \tag{5}
\]
Every $w\in\mathbb R^3$ is $Tu$ for some $u$, and adding a common
constant to the four entries of $u$ leaves $w$ unchanged. This makes
all $u_h$ positive. Pairwise distinctness of $u$ excludes only six
proper linear hyperplanes in $w$.

Put $z_a=1/d_a$. Multiplication by the positive number $d_a$ reduces
the desired limiting signs to those of
\[
 A_\ell+B_\ell z_a+w_\ell(\gamma z_a^2+\delta z_a^3).
 \tag{6}
\]
We prove that a function of this form realizes any three prescribed
sign columns on four distinct nodes.

Let $K\subset\mathbb R^4$ be the two-dimensional kernel of the
matrix with rows $(1,1,1,1)$ and $(z_1,z_2,z_3,z_4)$. Each coordinate
hyperplane cuts $K$ in a line, and these four lines are distinct.
For otherwise a nonzero vector supported on at most two coordinates
would belong to the span of $1,z$; an affine function vanishing at
the other two distinct nodes must vanish identically. The four lines
therefore divide $K$ into eight open sectors, yielding four antipodal
pairs of strict sign patterns. Among those four pairs, at least one
is absent from the three requested sign columns. Choose a vector
$q\in K$ in that sector. All entries of $q$ are nonzero.

Choose $p\in q^\perp$ outside $\operatorname{span}\{1,z\}$.
The four vectors of values of $1,z,p$ have rank three and left
kernel $\mathbb Rq$. For a requested sign column $f$, the numbers
$q_af_a$ contain both signs by the choice of $q$. Thus there are
strictly positive $m_a$ with $\sum_aq_af_am_a=0$. Explicitly, let
$P=\sum_{q_af_a>0}q_af_a$ and $N=-\sum_{q_af_a<0}q_af_a$, and take
$m_a=N$ on the positive terms and $m_a=P$ on the negative terms.
The vector $(f_am_a)_a$ belongs to $q^\perp$, so it is a linear
combination of $1,z,p$. Doing this independently for the three
columns yields
\[
 f_{a\ell}(A'_\ell+B'_\ell z_a+w_\ell p_a)>0.
 \tag{7}
\]
Interpolate $p$ by a polynomial of degree at most three. Its constant
and linear terms can be absorbed into $A',B'$, giving (6).
The strict inequalities in (7) remain true after a sufficiently
small arbitrary perturbation of the vector $w$. Choose such a
perturbation outside the six exceptional hyperplanes above, solve
$Tu=w$, add a common constant to make $u>0$, and set $t_h=\sqrt{u_h}$.
Equations (3)--(5) then supply coefficients in (2) with a strict
limiting sign at every one of the twelve points. Finiteness gives
one common positive upper bound on $\epsilon$.

For rational $d_a$, choose rational $q,p$ and the coefficients in
(7); the same open conditions allow rational perturbations. If the
square roots $t_h$ are irrational, approximate them by distinct
positive rationals. The limiting functions in (5) depend continuously
on $t$, so a sufficiently close approximation preserves all strict
signs. The column operations in (3)--(4) then have rational
coefficients for every positive rational $\epsilon$.
\end{proof}

\begin{lemma}[The single common cubic lift stops at four columns]
For fixed distinct nodes $z_1,\ldots,z_4$, there exist four sign
columns which cannot all be represented as $A_\ell+B_\ell z_a+w_\ell p_a$
with one common vector $p$.
\end{lemma}
\begin{proof}
Choose one sign pattern from each of the four antipodal sectors of
$K=\ker(1,z)^\top$. If $1,z,p$ has rank three, its left kernel is a
line in $K$. If that line has no zero coordinate, one of the four
chosen patterns, up to overall sign, is its strict sign pattern and
is impossible. If it has a zero coordinate, approach its line from
either adjacent sector. At least one chosen pattern agrees with the
nonzero signs of this left-kernel vector, up to overall sign, and
is again impossible: its dot product with any vector in that strict
orthant has a fixed nonzero sign. If the rank is only two, its image
is contained in the rank-three image obtained by adding any $p'$
outside $\operatorname{span}\{1,z\}$, so the same impossibility holds.
\end{proof}

\paragraph{Exact original-model conversion and its cost.}
The theorem is about masked columns, and those masks cannot be
silently included in an affine numerator. Here is an explicit
conversion with its full cost. Fix any finite number $s$ of memory
columns, choose the four $d_a$ distinct and sufficiently close to
one, and apply the theorem independently to each column $j$.
Index the twelve rows by $i=(a,\ell)$. For each of the eight signed
slots $h$, set its Boolean row mask to $B_{ih}=(1+V_{ih})/2$, where
$V_{ih}$ is the corresponding entry of $v_h$ or $-v_h$. This mask
matrix has rank at most three since it has only three distinct rows.

The unmasked denominator at $(i,j)$ is $d_a-\epsilon\beta_{hj}$,
with $\beta_{hj}=t_{hj}$ or $-t_{hj}$. Choose a number $C$ exceeding
all $|\beta_{hj}|$, an intercept $b>\max_a d_a$, and set
\[
 p_{hi}=b-d_a-\epsilon C,\qquad
 u_{hj}=\epsilon(C+\beta_{hj}).
 \tag{8}
\]
Then $b-p_{hi}-u_{hj}=d_a-\epsilon\beta_{hj}$. By taking the $d_a$
close enough to one, then $b$ slightly above their maximum, and
finally $\epsilon$ sufficiently small, all $p,u$ are positive and
$b-\sum_i p_{hi}-\sum_j u_{hj}>0$. Thus every denominator is
positive on the whole Boolean cube, with strictly negative slopes.
The unmasked head has column numerator $q_{hj}$, which is affine on
the one-hot column grid and independent of the row.

For completeness, a binary mask matrix $B$ of rank at most $L$ can
be applied to $H$ such heads using $H+L+1$ heads. In head $h$, replace
$p_{hi}$ by $B_{ih}p_{hi}$. Removing negative slopes preserves
whole-cube positivity. The new grid value is
$B_{ih}R_h(i,j)+(1-B_{ih})q_{hj}/(b-u_{hj})$. After factoring $B$ as
$B_{ih}=\sum_{\ell=1}^L a_{i\ell}b_{\ell h}$, the required correction
is a column term plus $L$ rank-one row-times-column terms.
Each rank-one term $F_iG_j$ costs one weak reverse head plus a column
term: choose $M$ large, use denominator $1/(M+F_i)$ on row $i$ and
numerator $G_j$, and subtract $MG_j$. This denominator has a positive
whole-cube extension for large $M$, since its row differences are
$O(M^{-2})$ while its intercept is $\Theta(M^{-1})$. Collect all
column terms into one constant-denominator head. This proves the
claimed count.

Since $V=2B-1$, subtracting the $H$ unmasked heads gives a total
$2H+L+1=20$. These may initially have zero denominator slopes. The
weight-two gauge adds $\tau(2-\sum_i x_i-\sum_j y_j)$ to each
denominator. It vanishes on the one-hot-by-one-hot grid, makes every
slope strictly negative, and preserves whole-cube positivity for
sufficiently small $\tau>0$. Finally any affine numerator over
$b-\sum_k a_kx_k$ with $a_k>0$ and $b>\sum_k a_k$ is an original
\texttt{FracAtom}: choose $0<\alpha<1$ with
$\sum_k a_k/(1-\alpha)<b$, set $\rho_k=a_k/(1-\alpha)$,
$\gamma=b-\sum_k\rho_k>0$, set $\delta=0$, and solve the numerator
coefficients for $m_k,\eta$. This checks the original admissibility
conditions on every Boolean input.

\paragraph{Audit.}
\texttt{audit\_reflection\_cauchy\_lift.py} uses exact rational
arithmetic at $z=1,2,3,4$. It verifies all 4096 sign tables, selects
an unused Gale sector, constructs the strict scores, and produces
positive distinct square tokens satisfying $Tu=w$. The four sector
choices occur 2744, 1016, 288, and 48 times. The script checks finite
limiting witnesses; it is not a Lean proof and does not assert an
asymptotic saving.
\end{document}
